\setcounter{section}{29}

  \zwischenueberschrift{Kelengkungan pada manifold Riemann}

Misalkan $M$ suatu
\definitionsverweis {manifold Riemann}{}{}
dan
\definitionsverweis {bundel tangen}{}{}
$TM$ dilengkapi dengan
\definitionsverweis {koneksi Levi--Civita}{}{.}
Kita hendak memahami
\definitionsverweis {operator kelengkungan}{}{}
$R$ dalam situasi ini secara lebih mendalam. Untuk medan-medan vektor $V,W$ yang diberikan, operator ini merupakan suatu pemetaan
\maabbdisp {R(V,W)} { C^2(M,TM) } { C^0(M,TM)
} {,}
yakni, ketika diterapkan pada suatu medan vektor $Z$
\zusatzklammer {yang dua kali terdiferensialkan} {} {,}
operator ini menghasilkan medan vektor

\mavergleichskettedisp
{\vergleichskette
{ R(V,W)(Z)
}
{ =} { \nabla_V \nabla_W Z
}
{ } {
}
{ } {
}
{ } {
}
}
{}{}{.}
Dengan sudut pandang ini, pada suatu manifold yang terdiferensialkan tak hingga, keseluruhan konstruksi merupakan pemetaan
\maabbdisp {} {C^\infty(TM) \times C^\infty(TM) \times C^\infty(TM)  } { C^\infty(TM)
} {,}
tetapi ketiga argumennya tidak mempunyai kedudukan yang setara.



\inputfaktbeweis
{R^n/Offene Menge/Riemannsche Metrik/Levi-Civita-Zusammenhang/Krümmungsoperator/Berechnung/Fakt}
{Lemma}
{}
{

\faktsituation {Misalkan

\mavergleichskette
{\vergleichskette
{ U
}
{ \subseteq }{ \R^n
}
{  }{
}
{    }{
}
{   }{
}
}
{}{}{}
\definitionsverweis {terbuka}{}{}
dan dilengkapi dengan suatu
\definitionsverweis {metrik Riemann}{}{}
$g_{ij}$. Misalkan $\nabla$ adalah
\definitionsverweis {koneksi Levi--Civita}{}{}
pada
\definitionsverweis {bundel tangen}{}{}

\mavergleichskette
{\vergleichskette
{ E
}
{ = }{ U \times \R^n
}
{  }{
}
{    }{
}
{   }{
}
}
{}{}{.}}
\faktfolgerung {Maka untuk
\definitionsverweis {operator kelengkungan}{}{}
berlaku

\mavergleichskettedisp
{\vergleichskette
{ R( \partial_i , \partial_j)( \partial_k)
}
{ =} { \sum_{\ell = 1}^n  R^\ell_{ijk} \partial_\ell
}
{ } {
}
{ } {
}
{ } {
}
}
{}{}{}
dengan

\mavergleichskettedisp
{\vergleichskette
{ R^\ell_{ijk}
}
{ =} { {\partial_i } \Gamma^\ell_{jk} - {\partial_j } \Gamma^\ell_{ik}  +  \sum_{ a  = 1}^n  \Gamma^a_{jk} \Gamma_{i a}^\ell -\sum_{ a  = 1}^n   \Gamma^a_{ik}  \Gamma_{j a}^\ell
}
{ } {
}
{ } {
}
{ } {
}
}
{}{}{,}
di mana

\mavergleichskettedisp
{\vergleichskette
{ \Gamma_{ij}^\ell
}
{ =} { { \frac{   1  }{  2 } } \sum_{\mu = 1}^n g^{\ell  \mu }(\partial_ig_{j\mu }+\partial_jg_{i\mu }-\partial_\mu g_{ij})
}
{ } {
}
{ } {
}
{ } {
}
}
{}{}{}
menyatakan
\definitionsverweis {simbol-simbol Christoffel}{}{}
dari koneksi tersebut.}
\faktzusatz {}
\faktzusatz {}

}
{

Ini merupakan kasus khusus dari
Lemma 28.2.

}



\inputfaktbeweis
{Riemannsche Mannigfaltigkeit/Levi-Civita-Zusammenhang/Krümmungsoperator/Eigenschaften/Fakt}
{Lemma}
{}
{

\faktsituation {Misalkan $M$ suatu
\definitionsverweis {manifold Riemann}{}{}
dan
\definitionsverweis {bundel tangen}{}{}
dilengkapi dengan
\definitionsverweis {koneksi Levi--Civita}{}{.}}
\faktuebergang {Maka
\definitionsverweis {operator kelengkungan}{}{}
mempunyai sifat-sifat berikut.}
\faktfolgerung {\aufzaehlungsechs{Ekspresi
\mathl{R(V,W)Z}{} linear dalam ketiga komponennya.
}{Untuk

\mavergleichskette
{\vergleichskette
{ P
}
{ \in }{ M
}
{  }{
}
{    }{
}
{   }{
}
}
{}{}{}
nilai
\mathl{(R(V,W)Z) (P)}{} hanya bergantung pada
\mathl{V(P),W(P),Z(P)}{}.
}{Berlaku

\mavergleichskettedisp
{\vergleichskette
{ R(V,W)
}
{ =} { -R(W,V)
}
{ } {
}
{ } {
}
{ } {
}
}
{}{}{}
}{Berlaku

\mavergleichskettedisp
{\vergleichskette
{ R(V,W)Z+ R(W,Z)V+R(Z,V)W
}
{ =} { 0
}
{ } {
}
{ } {
}
{ } {
}
}
{}{}{.}
}{Berlaku

\mavergleichskettedisp
{\vergleichskette
{ \left\langle R(V,W) T  ,  Z   \right\rangle
}
{ =} { - \left\langle R(V,W) Z  ,  T   \right\rangle
}
{ } {
}
{ } {
}
{ } {
}
}
{}{}{.}
}{Berlaku

\mavergleichskettedisp
{\vergleichskette
{ \left\langle R(V,W) T  ,  Z   \right\rangle
}
{ =} { \left\langle R(T,Z) V  ,  W   \right\rangle
}
{ } {
}
{ } {
}
{ } {
}
}
{}{}{.}
}}
\faktzusatz {}
\faktzusatz {}

}
{

\aufzaehlungsechs{Linearitas dalam $Z$ berasal dari linearitas koneksi. Linearitas dalam $V$ dan $W$ mengikuti dari
Lemma 24.10
dan
Lemma 11.8.
}{Untuk ketergantungan terhadap $V$ dan $W$, pernyataan tersebut mengikuti dari
Lemma 28.4.
Untuk menunjukkan bahwa ketergantungan terhadap $Z$ juga hanya bergantung pada $Z(P)$, kita dapat bekerja dalam situasi lokal pada

\mavergleichskette
{\vergleichskette
{ U
}
{ \subseteq }{ \R^n
}
{  }{
}
{    }{
}
{   }{
}
}
{}{}{}
dan mengambil

\mavergleichskette
{\vergleichskette
{ V
}
{ = }{ \partial_i
}
{  }{
}
{    }{
}
{   }{
}
}
{}{}{,}

\mavergleichskette
{\vergleichskette
{ W
}
{ = }{ \partial_j
}
{  }{
}
{    }{
}
{   }{
}
}
{}{}{}
dan

\mavergleichskette
{\vergleichskette
{ Z
}
{ = }{ h \partial_k
}
{  }{
}
{    }{
}
{   }{
}
}
{}{}{}
dengan suatu fungsi
\maabb {h} { U } {\R
} {}
yang dua kali terdiferensialkan secara kontinu. Menurut
Lemma 25.10,
Teorema 25.5,
dan
Teorema Schwarz,
kita memperoleh

\mavergleichskettealigndrucklinks
{\vergleichskettealigndrucklinks
{ R \left(  \partial_i, \partial_j  \right) \left(  h \partial_k  \right)
}
{ =} { \nabla_{\partial_i }  \left(  \nabla_{\partial_j }  \left(  h \partial_k  \right)  \right) - \nabla_{\partial_j } \left(  \nabla_{\partial_i }  \left(  h \partial_k  \right)  \right)
}
{ =} { \nabla_{\partial_i }  \left(  h \nabla_{\partial_j } \partial_k+ \left(  \partial_j h  \right)  \partial_k  \right) - \nabla_{\partial_j } \left( h \nabla_{\partial_i }  \partial_k  + \left(  \partial_i h  \right) \partial_k  \right)
}
{ =} { h \nabla_{\partial_i } \left(  \nabla_{\partial_j }  \partial_k   \right) +  \left(  \partial_i h  \right) \nabla_{\partial_j } \partial_k + \left(  \partial_j h  \right) \nabla_{\partial_i } \partial_k  +  \left(  \partial_i \partial_j h  \right) \partial_k
-h \nabla_{\partial_j } \left(  \nabla_{\partial_i } \partial_k   \right) -  \left(  \partial_j h  \right) \nabla_{\partial_i } \partial_k - \left(  \partial_i h  \right) \nabla_{\partial_j } \partial_k  -\left(  \partial_j  \partial_i h  \right) \partial_k
}
{ =} { h \nabla_{\partial_i } \left(  \nabla_{\partial_j }  \partial_k   \right)-h \nabla_{\partial_j } \left(  \nabla_{\partial_i } \partial_k   \right)
}
}
{}{}{,}
yang menunjukkan bahwa ekspresi ini hanya bergantung pada $h(P)$.
}{Hal ini langsung terlihat dari definisi dan
Lemma 11.8.
}{Karena
\definitionsverweis {koneksi bebas torsi}{}{}
\zusatzklammer {lihat
Teorema 26.14} {} {}
dan
identitas Jacobi,
berlaku

\mavergleichskettealigndrucklinks
{\vergleichskettealigndrucklinks
{ R(V,W)Z+ R(W,Z)V+R(Z,V)W
}
{ =} { \nabla_V  \left(  \nabla_W Z  \right) -\nabla_W  \left(  \nabla_V Z  \right) - \nabla_{[V,W]} Z +\nabla_W  \left(  \nabla_Z V  \right) -\nabla_Z  \left(  \nabla_W V  \right) - \nabla_{[W,Z]} V+\nabla_Z  \left(  \nabla_V W  \right) -\nabla_V  \left(  \nabla_Z W  \right) - \nabla_{[Z,V]} W
}
{ =} { \nabla_V  \left(  [W,Z ]  \right) + \nabla_W  \left(  [Z, V]  \right) + \nabla_Z  \left(  [V, W]  \right) - \nabla_{[V,W]} Z - \nabla_{[W,Z]} V  - \nabla_{[Z,V]} W
}
{ =} {\nabla_V  \left(  [W,Z ]  \right)  - \nabla_{[W,Z]} V + \nabla_W  \left(  [Z, V]  \right) - \nabla_{[Z,V]} W + \nabla_Z  \left(  [V, W]  \right) - \nabla_{[V,W]} Z
}
{ =} { [V, [W,Z]] + [W, [Z,V]] + [Z, [V,W]]
}
}
{
\vergleichskettefortsetzungalign
{ =} { 0
}
{ } {}
{ } {}
{ } {}
}{}{.}
}{Kita dapat membatasi diri pada medan-medan vektor $V,W$ dengan

\mavergleichskette
{\vergleichskette
{ [V,W]
}
{ = }{ 0
}
{  }{
}
{    }{
}
{   }{
}
}
{}{}{.}
Menurut
Teorema 26.14 (2)
dan
Teorema Schwarz,
kita kemudian memperoleh

\mavergleichskettealigndrucklinks
{\vergleichskettealigndrucklinks
{ \left\langle R(V,W)T , Z  \right\rangle
}
{ =} { \left\langle  \nabla_V \nabla_WT -\nabla_W \nabla_VT   ,  Z  \right\rangle
}
{ =} { \left\langle  \nabla_V \nabla_WT , Z  \right\rangle - \left\langle  \nabla_W \nabla_VT  ,  Z  \right\rangle
}
{ =} { - \left\langle  \nabla_WT , \nabla_V Z  \right\rangle + D_V\left\langle  \nabla_W T  ,  Z  \right\rangle + \left\langle  \nabla_VT  , \nabla_W  Z  \right\rangle -  D_W \left\langle  \nabla_V T  ,  Z  \right\rangle
}
{ =} { \left\langle  T ,  \nabla_W  \nabla_V Z  \right\rangle - D_W \left\langle  T ,  \nabla_VZ   \right\rangle  + D_V \left(  - \left\langle  T  ,  \nabla_WZ  \right\rangle +  D_W \left\langle T  ,  Z   \right\rangle  \right)  - \left\langle  T  ,  \nabla_V \nabla_W  Z  \right\rangle +  D_V \left\langle  T  ,  \nabla_W Z  \right\rangle - D_W  \left(  - \left\langle T  ,  \nabla_V Z   \right\rangle + D_V \left\langle T  ,  Z   \right\rangle  \right)
}
}
{
\vergleichskettefortsetzungalign
{ =} { \left\langle  T ,  \nabla_W  \nabla_V  Z  \right\rangle - \left\langle  T  ,  \nabla_V \nabla_W  Z  \right\rangle
}
{ =} { - \left\langle R(V,W)Z , T  \right\rangle
}
{ } {}
{ } {}
}{}{.}
}{Dengan menggunakan (3), (4), dan (5), diperoleh

\mavergleichskettealigndrucklinks
{\vergleichskettealigndrucklinks
{ \left\langle R(V,W)T , Z  \right\rangle
}
{ =} { - \left\langle R(T,V)W , Z  \right\rangle - \left\langle R(W,T)V , Z  \right\rangle
}
{ =} { \left\langle R(T,V) Z , W  \right\rangle + \left\langle R(W,T) Z ,  V  \right\rangle
}
{ =} { - \left\langle R(V,Z) T , W  \right\rangle - \left\langle R(Z,T) V  ,  W  \right\rangle - \left\langle R(T,Z) W ,  V  \right\rangle - \left\langle R(Z,W) T ,  V  \right\rangle
}
{ =} { 2 \left\langle R(T,Z) V  ,  W  \right\rangle - \left\langle R(V,Z) T , W  \right\rangle - \left\langle R(Z,W) T ,  V  \right\rangle
}
}
{
\vergleichskettefortsetzungalign
{ =} { 2 \left\langle R(T,Z) V  ,  W  \right\rangle + \left\langle R(V,Z) W , T  \right\rangle + \left\langle R(Z,W) V  ,  T  \right\rangle
}
{ =} { 2 \left\langle R(T,Z) V  ,  W  \right\rangle - \left\langle R(W,V) Z , T  \right\rangle
}
{ =} { 2 \left\langle R(T,Z) V  ,  W  \right\rangle - \left\langle R(V,W) T , Z  \right\rangle
}
{ } {}
}{}{.}
Ini membuktikan pernyataan tersebut.
}

}

Berdasarkan
Lemma 29.2 (2),
untuk vektor-vektor

\mavergleichskette
{\vergleichskette
{ v,w,z
}
{ \in }{ T_PM
}
{  }{
}
{    }{
}
{   }{
}
}
{}{}{}
di suatu titik

\mavergleichskette
{\vergleichskette
{ P
}
{ \in }{ M
}
{  }{
}
{    }{
}
{   }{
}
}
{}{}{}
terdefinisi dengan baik suatu vektor

\mavergleichskette
{\vergleichskette
{ R(v,w)(z)
}
{ \in }{ T_PM
}
{  }{
}
{    }{
}
{   }{
}
}
{}{}{,}
sebab vektor-vektor tersebut dapat direalisasikan oleh medan-medan vektor terdiferensialkan
\mathl{V,W,Z}{} dengan

\mavergleichskette
{\vergleichskette
{ V(P)
}
{ = }{ v
}
{  }{
}
{    }{
}
{   }{
}
}
{}{}{}
dan seterusnya
\zusatzklammer {secara lokal} {} {}
dan kemudian $R(V,W)(Z)$ dapat dievaluasi di titik tersebut.

  \zwischenueberschrift{Kelengkungan seksional}

\inputdefinition
{}
{

Misalkan $M$ suatu
\definitionsverweis {manifold Riemann}{}{,}
yang dilengkapi dengan
\definitionsverweis {koneksi Levi--Civita}{}{}
pada
\definitionsverweis {bundel tangen}{}{}
$TM$ dan
\definitionsverweis {operator kelengkungan}{}{}
terkait $R$. Untuk vektor-vektor
\definitionsverweis {bebas linear}{}{}

\mavergleichskette
{\vergleichskette
{ v,w
}
{ \in }{ T_PM
}
{  }{
}
{    }{
}
{   }{
}
}
{}{}{}
di atas suatu titik

\mavergleichskette
{\vergleichskette
{ P
}
{ \in }{ M
}
{  }{
}
{    }{
}
{   }{
}
}
{}{}{}
bilangan

\mavergleichskettedisp
{\vergleichskette
{ K(v,w)
}
{ =} { { \frac{   \left\langle R(v,w)w , v  \right\rangle  }{  \left\langle  v  , v  \right\rangle  \left\langle  w  , w  \right\rangle  -\left\langle  v  , w  \right\rangle^2   } }
}
{ } {
}
{ } {
}
{ } {
}
}
{}{}{}
disebut
\definitionswort {kelengkungan seksional}{}
untuk $v,w$ di $P$.

}

Kelengkungan seksional ini terdefinisi dengan baik karena, dalam kasus kebebasan linear, penyebutnya tidak sama dengan $0$.



\inputfaktbeweis
{Riemannsche Mannigfaltigkeit/Levi-Civita-Zusammenhang/Schnittkrümmung/Tangentiale Ebene/Fakt}
{Lemma}
{}
{

\faktsituation {Misalkan $M$ suatu
\definitionsverweis {manifold Riemann}{}{,}
yang dilengkapi dengan
\definitionsverweis {koneksi Levi--Civita}{}{}
pada
\definitionsverweis {bundel tangen}{}{}
$TM$ dan
\definitionsverweis {operator kelengkungan}{}{}
terkait $R$.}
\faktfolgerung {Maka
\definitionsverweis {kelengkungan seksional}{}{}
untuk dua vektor bebas linear

\mavergleichskette
{\vergleichskette
{ v,w
}
{ \in }{ T_PM
}
{  }{
}
{    }{
}
{   }{
}
}
{}{}{}
hanya bergantung pada bidang yang direntang oleh vektor-vektor tersebut,

\mavergleichskettedisp
{\vergleichskette
{ \R v + \R w
}
{ \subseteq} { T_PM
}
{ } {
}
{ } {
}
{ } {
}
}
{}{}{.}}
\faktzusatz {}
\faktzusatz {}

}
{

Kita meninjau ekspresi pendefinisi

\mavergleichskettedisp
{\vergleichskette
{ K(v,w)
}
{ =} { { \frac{   \left\langle R(v,w)w , v  \right\rangle  }{  \left\langle  v  , v  \right\rangle \left\langle  w  , w  \right\rangle  - \left\langle  v  , w  \right\rangle^2   } }
}
{ } {
}
{ } {
}
{ } {
}
}
{}{}{}
untuk kelengkungan seksional. Menurut
Lemma 29.2,
ekspresi
\mathl{R(v,w)z}{} linear dalam semua komponennya. Jika $v$ atau $w$ dikalikan dengan suatu skalar

\mavergleichskette
{\vergleichskette
{ a
}
{ \neq }{ 0
}
{  }{
}
{    }{
}
{   }{
}
}
{}{}{,}
faktor $a^2$ dapat dikeluarkan baik dari pembilang maupun penyebut. Selanjutnya,

\mavergleichskettealignhandlinks
{\vergleichskettealignhandlinks
{ \left\langle  R(v+w,w) w  , v +w   \right\rangle
}
{ =} { \left\langle  R(v,w)w +R(w,w)w  , v+w   \right\rangle
}
{ =} { \left\langle  R(v,w)w   , v   \right\rangle  +  \left\langle  R(w,w)w  ,  v   \right\rangle + \left\langle  R(v,w)w  ,  w   \right\rangle  +  \left\langle  R(w,w)w   , w   \right\rangle
}
{ } {
}
{ } {
}
}
{}
{}{.}
Menurut
Lemma 29.2 (4),
\mathkor {} {\left\langle  R(w,w)w  ,  v   \right\rangle} {dan} {\left\langle  R(w,w)w   , w   \right\rangle} {}
sama dengan $0$, dan menurut
Lemma 29.2 (5,6),
juga berlaku

\mavergleichskette
{\vergleichskette
{ \left\langle  R(v,w)w  ,  w   \right\rangle
}
{ = }{ 0
}
{  }{
}
{    }{
}
{   }{
}
}
{}{}{.}
Jadi pembilang kelengkungan seksional tidak berubah ketika $v$ diganti dengan $v+w$. Karena

\mavergleichskettealigndrucklinks
{\vergleichskettealigndrucklinks
{ \left\langle v+w , v+w  \right\rangle \left\langle  w  , w  \right\rangle -\left\langle v+w , w  \right\rangle^2
}
{ =} { \left\langle  v  , v  \right\rangle  \left\langle  w  , w  \right\rangle+\left\langle  w  , w  \right\rangle \left\langle  w  , w  \right\rangle +2\left\langle  v  , w  \right\rangle \left\langle  w  , w  \right\rangle -\left\langle v  , w  \right\rangle^2-\left\langle w  , w  \right\rangle^2 -2 \left\langle v  , w  \right\rangle \left\langle w  , w  \right\rangle
}
{ =} { \left\langle  v  , v  \right\rangle \left\langle  w  , w  \right\rangle -\left\langle v  , w  \right\rangle^2
}
{ } {
}
{ } {
}
}
{}{}{}
penyebutnya juga tidak berubah.

}

\inputbemerkung
{}
{

Pada suatu manifold Riemann berdimensi dua, satu-satunya subruang vektor berdimensi dua dalam setiap ruang tangen adalah ruang tangen itu sendiri. Ini berarti bahwa, dalam kasus berdimensi dua, menurut
Lemma 29.4,
argumen-argumen dalam
\definitionsverweis {kelengkungan seksional}{}{}
tidak diperlukan. Karena itu, dalam situasi ini kita memandang kelengkungan seksional sebagai suatu fungsi bernilai riil pada manifold tersebut.

}



\inputfaktbeweis
{Riemannsche Mannigfaltigkeit/Zweidimensional/Levi-Civita-Zusammenhang/Schnittkrümmung/Berechnung/Fakt}
{Lemma}
{}
{

\faktsituation {Untuk suatu
\definitionsverweis {manifold Riemann}{}{}
berdimensi dua}
\faktfolgerung {
\definitionsverweis {kelengkungan seksional}{}{}
pada setiap peta sama dengan
\mathl{-  { \frac{   R^2_{121} }{ g_{11}  } }}{.}}
\faktzusatz {}
\faktzusatz {}

}
{

Kita dapat langsung meninjau situasi pada suatu peta. Menurut
Lemma 29.4,
dalam kasus berdimensi dua kelengkungan seksional dapat dihitung menggunakan basis sembarang; kita memilih
\mathkor {} {\partial_1} {dan} {\partial_2} {.}
Dengan menggunakan
Lemma 29.2 (5),
untuk kelengkungan seksional diperoleh

\mavergleichskettealign
{\vergleichskettealign
{ K( \partial_1, \partial_2)
}
{ =} { { \frac{   \left\langle R(\partial_1 ,\partial_2)\partial_2  ,  \partial_1   \right\rangle  }{  \left\langle  \partial_1  ,  \partial_1   \right\rangle  \left\langle  \partial_2  ,  \partial_2   \right\rangle  -\left\langle  \partial_1  ,  \partial_2   \right\rangle^2   } }
}
{ =} { { \frac{   \left\langle R(\partial_1 ,\partial_2)\partial_2  ,  \partial_1   \right\rangle  }{  g_{11} g_{22}  -g_{12}^2   } }
}
{ =} { -  { \frac{   \left\langle R(\partial_1 ,\partial_2)\partial_1  ,  \partial_2   \right\rangle  }{  g_{11} g_{22}  -g_{12}^2   } }
}
{ } {
}
}
{}
{}{}
Menurut
Lemma 29.1,
berlaku

\mavergleichskettedisp
{\vergleichskette
{ R( \partial_1 , \partial_2)( \partial_1)
}
{ =} { R^1_{121} \partial_1+ R^2_{121} \partial_2
}
{ } {
}
{ } {
}
{ } {
}
}
{}{}{.}
Hal ini memberikan

\mavergleichskettealign
{\vergleichskettealign
{ K( \partial_1, \partial_2)
}
{ =} { - { \frac{   \left\langle R(\partial_1 ,\partial_2)\partial_1  ,  \partial_2   \right\rangle  }{  g_{11} g_{22}  -g_{12}^2   } }
}
{ =} { - { \frac{   R^1_{121} g_{12}+ R^2_{121} g_{22}  }{  g_{11} g_{22}  -g_{12}^2   } }
}
{ } {
}
{ } {
}
}
{}
{}{.}
Menurut
Lemma 29.2 (4),

\mavergleichskettedisp
{\vergleichskette
{ \left\langle  R(\partial_1, \partial_1) \partial_1  ,  \partial_2   \right\rangle
}
{ =} { R^1_{121}  g_{11}  + R_{121}^2 g_{12}
}
{ =} { 0
}
{ } {
}
{ } {
}
}
{}{}{.}
Kita mengalikan pembilang kelengkungan seksional dengan $g_{11}$ dan memperoleh

\mavergleichskettealignhandlinks
{\vergleichskettealignhandlinks
{ g_{11}  \left(  R^1_{121} g_{12}+ R^2_{121} g_{22}  \right)
}
{ =} { g_{11}  \left(  R^1_{121} g_{12}+ R^2_{121} g_{22}  \right) - g_{12}  \left(  R^1_{121}  g_{11}  + R_{121}^2 g_{12}   \right)
}
{ =} { g_{11} g_{22} R^2_{121} - g_{12}^2 R^2_{121}
}
{ =} { \left(  g_{11} g_{22}  - g_{12}^2  \right) R^2_{121}
}
{ } {
}
}
{}
{}{.}
Jadi,

\mavergleichskettedisp
{\vergleichskette
{ K( \partial_1, \partial_2)
}
{ =} { -  { \frac{   R^2_{121} }{ g_{11}  } }
}
{ } {
}
{ } {
}
{ } {
}
}
{}{}{.}

}



\inputfaktbeweis
{Fläche im Raum/Schnittkrümmung/Gaußkrümmung/Fakt}
{Lemma}
{}
{

\faktsituation {Misalkan

\mavergleichskette
{\vergleichskette
{ W
}
{ \subseteq }{ \R^3
}
{  }{
}
{    }{
}
{   }{
}
}
{}{}{}
\definitionsverweis {terbuka}{}{}
dan misalkan
\maabb {h} { W } { \R
} {}
suatu
\definitionsverweis {fungsi yang dua kali terdiferensialkan secara kontinu}{}{}
serta

\mavergleichskette
{\vergleichskette
{ Y
}
{ = }{ h^{-1}(c)
}
{  }{
}
{    }{
}
{   }{
}
}
{}{}{}
serat di atas

\mavergleichskette
{\vergleichskette
{ c
}
{ \in }{ \R
}
{  }{
}
{    }{
}
{   }{
}
}
{}{}{,}
dan andaikan $h$
\definitionsverweis {reguler}{}{}
di setiap titik $Y$.}
\faktfolgerung {Maka
\definitionsverweis {kelengkungan Gauss}{}{}
dari $Y$ sama dengan
\definitionsverweis {kelengkungan seksional}{}{}
dari $Y$.}
\faktzusatz {}
\faktzusatz {}

}
{

Menurut
Lemma 29.6,
kelengkungan seksional pada suatu peta sama dengan
\mathl{- { \frac{  R^2_{121}  }{ g_{11}  } }}{.} Karena

\mavergleichskettedisphandlinks
{\vergleichskettedisphandlinks
{ R^2_{121}
}
{ =} { {\partial_1 } \Gamma^2_{12} - {\partial_2 } \Gamma^2_{11} +  \Gamma^1_{12} \Gamma_{1 1}^2 + \Gamma^2_{12} \Gamma_{1 2}^2 - \Gamma^1_{11}  \Gamma_{12 }^2  - \Gamma^2_{11} \Gamma_{2 2}^2
}
{ =} { -g_{11} K
}
{ } {
}
{ } {
}
}
{}{}{.}
maka

\mavergleichskettedisphandlinks
{\vergleichskettedisphandlinks
{ K
}
{ =} { { \frac{   1  }{ g_{11}  } }  \left(  \partial_2 \Gamma_{11}^2 - \partial_1 \Gamma_{12}^2 + \Gamma^1_{11} \Gamma^2_{12}  + \Gamma_{11}^2 \Gamma_{22}^2 - \Gamma_{12}^1 \Gamma_{11}^2-  \left(  \Gamma_{12}^2  \right)^2  \right)
}
{ } {
}
{ } {
}
{ } {
}
}
{}{}{}
yang berimpit dengan kelengkungan Gauss menurut
Teorema 19.7.

}

\inputbeispiel{}
{

Untuk
\definitionsverweis {bidang Euklides}{}{}
$\R^2$ dengan
\definitionsverweis {hasil kali skalar standar}{}{,}
\definitionsverweis {kelengkungan seksional}{}{}
konstan sama dengan $0$.

}

\inputbeispiel{}
{

Untuk
\definitionsverweis {bola satuan}{}{}

\mavergleichskette
{\vergleichskette
{ S^2
}
{ \subseteq }{ \R^3
}
{  }{
}
{    }{
}
{   }{
}
}
{}{}{}
dengan
\definitionsverweis {struktur Riemann}{}{}
terinduksi,
\definitionsverweis {kelengkungan seksional}{}{}
konstan sama dengan $1$. Hal ini mengikuti dari
Contoh 5.6
bersama dengan
Lemma 29.7.

}

\inputbeispiel{}
{

\bild{ \begin{center}
\includegraphics[width=5.5cm]{ \bildeinlesungsvg {Parallel lines in Poincare's model of hyperbolic geometry} {svg} }
\end{center}
\bildtext {} }

\bildlizenz { Parallel lines in Poincare's model of hyperbolic geometry.svg } {} {Januszkaja} {Commons} {CC-by-sa 3.0} {}

Kita melanjutkan
Contoh 18.8
dan
Contoh 26.7.
Menurut
Lemma 29.1,

\mavergleichskettedisphandlinks
{\vergleichskettedisphandlinks
{ R^2_{121}
}
{ =} { {\partial_1 } \Gamma^2_{12} - {\partial_2 } \Gamma^2_{11} + \Gamma^1_{12} \Gamma_{1 1}^2 + \Gamma^2_{12} \Gamma_{1 2}^2 - \Gamma^1_{11} \Gamma_{12 }^2 - \Gamma^2_{11} \Gamma_{2 2}^2
}
{ =} { { \frac{   1  }{  y^2 } } - { \frac{   1  }{  y^2 } }+  { \frac{   1  }{  y^2 } }
}
{ =} { { \frac{   1  }{  y^2 } }
}
{ } {
}
}
{}{}{.}
Menurut
Lemma 29.6,
untuk
\definitionsverweis {kelengkungan seksional}{}{}
berlaku

\mavergleichskettedisp
{\vergleichskette
{ K
}
{ =} { - { \frac{   R^2_{121} }{ g_{11 } }}
}
{ =} { - { \frac{   { \frac{   1  }{  y^2 } }  }{  { \frac{   1  }{  y^2 } }  } }
}
{ =} { -1
}
{ } {
}
}
{}{}{,}
jadi kita memperoleh kelengkungan seksional konstan $-1$.

}
